\section{Omitted Proofs}

\subsection{Proof of Theorem~\ref{thm:pivot}}


\begin{proof}
Fix $s$ and $\beta>0$. Write for simplicity
\begin{align*}
\pi = \pi_{s,\beta}, \quad
q = \mathbb E_{a\sim\pi}[r(s,a)], \quad
\sigma^2 = \mathrm{Var}_{a\sim\pi}(r(s,a)).
\end{align*}

We first compute the natural gradient of $J_s$ at $\pi$.
Let $v\in T_\pi$, and consider the exponential perturbation
\begin{align*}
\pi_t(a)= &~ \frac{\pi(a)e^{t v(a)}}{\sum_{b\in\mathcal{A}(s) }\pi(b)e^{t v(b)}}.
\end{align*}
For any $v\in T_\pi$, we have
\begin{align*}
\left.\frac{d}{dt}\right|_{t=0}\log \pi_t(a)
= v(a)-\sum_{b\in\mathcal{A}(s) }\pi(b)v(b)
= v(a).
\end{align*}
Thus
\begin{align*}
DJ_s(\pi)[v]
= &~ \left.\frac{d}{dt}\right|_{t=0}\mathbb E_{a\sim\pi_t}[r(s,a)] \\
= &~ \left.\frac{d}{dt}\right|_{t=0}\sum_{a\in\mathcal{A}(s) }\pi_t(a)r(s,a) \\
= &~ \sum_{a\in\mathcal{A}(s) }\pi(a)r(s,a)v(a) \\
= &~ \sum_{a\in\mathcal{A}(s) }\pi(a)\bigl(r(s,a)-q\bigr)v(a) \\
= &~ \langle r(s,\cdot)-q,\,v\rangle_{F,\pi}.
\end{align*}
Since this holds for every $v\in T_\pi$, the Riesz representative is
\begin{align*}
\nabla^{\mathrm{nat}}J_s(\pi)= &~ r(s,\cdot)-q.
\end{align*}
We then have
\begin{align}\label{eq:norm-natural-J}
\bigl\|\nabla^{\mathrm{nat}}J_s(\pi)\bigr\|_{F,\pi}^2
= &~ \mathbb E_{a\sim\pi}\bigl[(r(s,a)-q)^2\bigr] \notag\\
= &~ \mathrm{Var}_{a\sim\pi}(r(s,a)) \notag\\
= &~ \sigma^2
\end{align}

Next, we compute the $\beta$-derivative of the exponential-tilted path. Define
\begin{align*}
Z_s(\beta)= &~ \sum_{b\in\mathcal{A}(s) }\pi_0(b\mid s)e^{r(s,b)/\beta}.
\end{align*}
Then
\begin{align*}
\log \pi_{s,\beta}(a)
= &~ \log \pi_0(a\mid s) + \frac{r(s,a)}{\beta} - \log Z_s(\beta).
\end{align*}
Differentiating with respect to $\beta$ gives
\begin{align*}
\partial_\beta \log \pi_{s,\beta}(a)
= &~ -\frac{r(s,a)}{\beta^2} - \frac{Z_s'(\beta)}{Z_s(\beta)}.
\end{align*}
Now
\begin{align*}
Z_s'(\beta)
= &~ \sum_{b\in\mathcal{A}(s) }\pi_0(b\mid s)e^{r(s,b)/\beta}\left(-\frac{r(s,b)}{\beta^2}\right) \\
= &~ -\frac{1}{\beta^2}\sum_{b\in\mathcal{A}(s) }\pi_0(b\mid s)e^{r(s,b)/\beta}r(s,b) \\
= &~ -\frac{Z_s(\beta)}{\beta^2}\,\mathbb E_{b\sim\pi_{s,\beta}}[r(s,b)] \\
= &~ -\frac{q}{\beta^2}Z_s(\beta).
\end{align*}
Therefore
\begin{align*}
\partial_\beta \log \pi_{s,\beta}(a)
= &~ -\frac{1}{\beta^2}\bigl(r(s,a)-q\bigr) \\
= &~ -\frac{1}{\beta^2}\nabla^{\mathrm{nat}}J_s(\pi_{s,\beta})(a).
\end{align*}

If $\sigma=0$, then $r(s,a)=q$ for $\pi$-almost every $a$, so
\begin{align*}
\nabla^{\mathrm{nat}}J_s(\pi)= &~ 0, ~~\text{and}~~ \frac{1}{\beta^2}\bigl\|\nabla^{\mathrm{nat}}J_s(\pi)\bigr\|_{F,\pi}= 0.
\end{align*}

Now assume $\sigma>0$. Then
\begin{align*}
\gamma_{s,\beta}
= &~ -\mathbb E_{a\sim\pi}
\left[
\frac{r(s,a)-q}{\sigma}\,\partial_\beta \log \pi(a)
\right] \\
= &~ -\mathbb E_{a\sim\pi}
\left[
\frac{r(s,a)-q}{\sigma}\left(-\frac{1}{\beta^2}(r(s,a)-q)\right)
\right] \\
= &~ \frac{1}{\beta^2\sigma}\mathbb E_{a\sim\pi}\bigl[(r(s,a)-q)^2\bigr] \\
= &~ \frac{\sigma}{\beta^2}.
\end{align*}
Combining this with Eq.~\eqref{eq:norm-natural-J}, we have
\begin{align*}
\gamma_{s,\beta}
= &~ \frac{1}{\beta^2}\bigl\|\nabla^{\mathrm{nat}}J_s(\pi_{s,\beta})\bigr\|_{F,\pi_{s,\beta}} = \frac{\sigma}{\beta^2}.
\end{align*}
This proves the theorem.
\end{proof}


\subsection{Proof of Theorem~\ref{thm:kl}}\label{sec:proof-kl-proj}

\begin{proof}
Since
\begin{align*}
\mathcal L_{\mathrm{func},\beta}(\pi)= &~ \mathbb E_{s\sim d}[\ell_s(\pi)]
\end{align*}
has no coupling across states, it suffices to solve the pointwise minimization problem for each fixed state $s$. The claim for the population objective then follows immediately $d$-almost surely.

Fix a state $s$, and write $M= \mathcal M(s), ~\rho= \pi_0(M\mid s), ~\pi_0(a)= \pi_0(a\mid s)$.

We first handle the boundary cases. 
If $\rho=0$, then $M=\varnothing$ because $\pi_0(a)>0$ for every $a\in \mathcal{A}(s) $. Hence $r_{\mathrm{func}}(s,a)\equiv 0$, so $\ell_s(\pi)= \beta\, \mathrm{KL}(\pi\|\pi_0)$,
which is uniquely minimized at $\pi=\pi_0$. 
If $\rho=1$, then $M=\mathcal{A}(s) $, so $r_{\mathrm{func}}(s,a)\equiv 1$, and $\ell_s(\pi)= -1 + \beta\, \mathrm{KL}(\pi\|\pi_0)$,
which is again uniquely minimized at $\pi=\pi_0$.

It remains to consider the nondegenerate case $0<\rho<1$.
Let $p$ be any distribution on $\mathcal{A}(s) $, and let
\begin{align*}
q= &~ p(M)=\sum_{a\in M} p(a).
\end{align*}
For $q\in(0,1)$, define the conditional distributions
\begin{align*}
p_+(a)= \frac{p(a)}{q}, \quad a\in M, \qquad
p_-(a)= \frac{p(a)}{1-q}, \quad a\notin M,
\end{align*}
and similarly
\begin{align*}
\pi_{0,+}(a)= \frac{\pi_0(a)}{\rho}, \quad a\in M, \qquad
\pi_{0,-}(a)= \frac{\pi_0(a)}{1-\rho}, \quad a\notin M.
\end{align*}

Using the chain rule for KL divergence, we obtain
\begin{align*}
\mathrm{KL}(p\|\pi_0)
= &~ \sum_{a\in M} p(a)\log\frac{p(a)}{\pi_0(a)}
   + \sum_{a\notin M} p(a)\log\frac{p(a)}{\pi_0(a)} \\
= &~ q\,\mathrm{KL}(p_+\|\pi_{0,+})
   + (1-q)\,\mathrm{KL}(p_-\|\pi_{0,-}) + q\log\frac{q}{\rho}
   + (1-q)\log\frac{1-q}{1-\rho}.
\end{align*}

Since $\mathbb E_{a\sim p}[r_{\mathrm{func}}(s,a)]=q$, the pointwise objective becomes
\begin{align*}
\ell_s(p)
= -q
+ \beta\Bigg[
q\,\mathrm{KL}(p_+\|\pi_{0,+})
+ (1-q)\,\mathrm{KL}(p_-\|\pi_{0,-})
+ q\log\frac{q}{\rho}
+ (1-q)\log\frac{1-q}{1-\rho}
\Bigg].
\end{align*}

For fixed $q$, the last two terms are constants, while the first two KL terms are nonnegative and vanish if and only if
\begin{align*}
p_+= \pi_{0,+}, \quad
p_-= \pi_{0,-}.
\end{align*}
Hence, for fixed $q$, the unique minimizer is
\begin{align*}
\bar p_q(a)= &~
\begin{cases}
\displaystyle \frac{q}{\rho}\,\pi_0(a), & a\in M,\\[1.1ex]
\displaystyle \frac{1-q}{1-\rho}\,\pi_0(a), & a\notin M.
\end{cases}
\end{align*}

Substituting $p=\bar p_q$ eliminates the conditional KL terms and yields the scalar problem
\begin{align*}
\phi_s(q)= &~ -q
+ \beta\left[
q\log\frac{q}{\rho}
+ (1-q)\log\frac{1-q}{1-\rho}
\right],
\qquad q\in[0,1],
\end{align*}
with the usual convention $0\log 0 = 0$.

For $q\in(0,1)$,
\begin{align*}
\phi_s'(q)= &~ -1 + \beta\log\frac{q(1-\rho)}{(1-q)\rho}, \\
\phi_s''(q)= &~ \beta\left(\frac{1}{q}+\frac{1}{1-q}\right) > 0.
\end{align*}
Therefore $\phi_s$ is strictly convex on $(0,1)$ and has at most one stationary point there. Solving $\phi_s'(q)=0$ gives $q= \frac{\rho e^{1/\beta}}{(1-\rho)+\rho e^{1/\beta}}$.
Hence the unique minimizer of $\ell_s$ is $\pi_\beta^\star(\cdot\mid s)=\bar p_{q_\beta(s)}$, namely
\begin{align}\label{eq:pi-beta-star}
\pi_\beta^\star(a\mid s)= &~
\begin{cases}
\displaystyle \frac{q_\beta(s)}{\rho(s)}\,\pi_0(a\mid s), & a\in \mathcal M(s),\\[1.1ex]
\displaystyle \frac{1-q_\beta(s)}{1-\rho(s)}\,\pi_0(a\mid s), & a\notin \mathcal M(s).
\end{cases}
\end{align}
Among all distributions with total acceptable mass $q_\beta(s)$, the unique KL minimizer is precisely the block-rescaled distribution above.

From Eq.~\eqref{eq:pi-beta-star}, the ordering-preservation statements are immediate. If $a,b\in \mathcal M(s)$, then
\begin{align*}
\frac{\pi_\beta^\star(a\mid s)}{\pi_\beta^\star(b\mid s)}
= &~ \frac{\pi_0(a\mid s)}{\pi_0(b\mid s)}.
\end{align*}
Likewise, if $a,b\notin \mathcal M(s)$, then
\begin{align*}
\frac{\pi_\beta^\star(a\mid s)}{\pi_\beta^\star(b\mid s)}
= &~ \frac{\pi_0(a\mid s)}{\pi_0(b\mid s)}.
\end{align*}
This completes the proof.
\end{proof}
